% ============================================================
%  訴状テンプレート【従来型・表紙あり・紙提出向け】(kian-v08 LuaLaTeX 版)
%  ※ 裁判IT化（電子提出）の場合は template-sojo-it-lua.tex を使用すること
%
%  コンパイル: latexmk -r latexmkrc-lua template-sojo-lua.tex
%              または VS Code（LaTeX Workshop）でビルド
%  プレースホルダー: 【...】を実際の内容に置き換えてください
%  異体字: \ivs{文字} で IPAmj 明朝字形に切り替え（例: \ivs{角󠄁}）
%  ※ \SecA{} 等の {} 内に % を書くとコメントになりエラーが出ます
%  見出し番号: \SecZ=番号なし \SecA=第１ \SecB=１ \SecC=⑴ \SecD=ア \SecE=(ア)
%  請求の趣旨も請求の原因と同じく \SecB{} で番号付けする（enumerate は使わない）。
% ============================================================

\documentclass{ltjsarticle}
\usepackage{kian-lua}
% \usepackage{kian-ivs} %IVS 異体字（\ivs{文字}）
% \usepackage{kian-otf} %囲み文字拡張
% \usepackage{luatexja-ruby} %ルビ機能

% ---------------------------------------------------------- %
%  表題・提出日・提出先の設定
% ---------------------------------------------------------- %

\Tit{訴　状}
\Dat{令和８年　　月　　日}
\Recs{
  [n]{:【裁判所名】【支部名】支部:御中}
}

% ============================================================

\begin{document}

\maketitle

% ---------------------------------------------------------- %
%  当事者欄
%  氏名に異体字がある場合: 【苗字】　\ivs{文字}　【名前】
% ---------------------------------------------------------- %

\noindent\textbf{原　告}

\noindent\Jus[4]{住　所}　【原告の住所】

\noindent\Jus[4]{氏　名}　【原告の氏名】

\vspace{.5\zw}

\noindent\textbf{上記訴訟代理人弁護士}

\noindent\Jus[4]{住　所}　【事務所住所】

\noindent\Jus[4]{氏　名}　【氏名】

\noindent\Jus[4]{電　話}　【電話番号】

\noindent\Jus[4]{ＦＡＸ}　【FAX番号】

\vspace{1\zw}

\noindent\textbf{被　告}

\noindent\Jus[4]{住　所}　【被告の住所】

\noindent\Jus[4]{氏　名}　【被告の氏名】

\vspace{1\zw}

% ---------------------------------------------------------- %
%  訴訟物の価額・貼用印紙額
% ---------------------------------------------------------- %

\noindent 訴訟物の価額　金【価額】円

\noindent 貼用印紙額　　金【印紙額】円

\vspace{1\zw}

% ---------------------------------------------------------- %
%  請求の趣旨
% ---------------------------------------------------------- %

\SecA{請求の趣旨}

\SecB{}

被告は原告に対し、金【請求金額】円及びこれに対する
令和【年】年【月】月【日】日から支払済みまで年【利率】の割合による
金員を支払え。

\SecB{}

訴訟費用は被告の負担とする。

% \SecB 後は字下げが１の本文位置になるため、第１の本文位置に戻す
{\advance\leftskip by -\parindent
との判決並びに仮執行宣言を求める。\par}

\vspace{1\zw}

% ---------------------------------------------------------- %
%  請求の原因
% ---------------------------------------------------------- %

\SecA{請求の原因}

\SecB{当事者}



\SecB{請求原因事実}



\SecB{損害}



\SecB{まとめ}



\vspace{1\zw}

% ---------------------------------------------------------- %
%  証拠方法
% ---------------------------------------------------------- %

\SecA{証拠方法}

\begin{enumerate}
  \item 甲第１号証　【証拠名称】
  \item 甲第２号証　【証拠名称】
  \item 甲第３号証　【証拠名称】
\end{enumerate}

\vspace{1\zw}

% ---------------------------------------------------------- %
%  附属書類
% ---------------------------------------------------------- %

\SecA{附属書類}

\begin{enumerate}
  \item 訴状副本　　　　　　　　　　　　　　　１通
  \item 甲号証写し　　　　　　　　　　　　　各１通
  \item 訴訟委任状　　　　　　　　　　　　　　１通
\end{enumerate}

\Ijo[2\zw]

\end{document}
